
\textbf{Corollaire 6.3.} \uline{Avec les hypothèses de} 6.1, \uline{soit} @@M0@@ \uline{un anneau commutatif}. \uline{Pour qu'un élément} @@M1@@ \uline{de} @@M2@@ \uline{soit inversible}, \uline{il faut et il suffit que} @@M3@@ \uline{le soit}. \uline{En particulier}, \uline{si} @@M4@@ \uline{est un faisceau localement libre de type fini sur} @@M5@@, \uline{partout} @@M6@@, \uline{alors sa classe dans} @@M7@@ \uline{est inversible}.

L'idéal d'augmentation de @@M8@@ est @@M9@@, et est donc un nilidéal, d'après 6.1. Pour qu'un élément de @@M10@@ soit inversible il est donc nécessaire et suffisant que son augmentation le soit.

\textbf{Corollaire 6.4.} \uline{Soient} @@M11@@ \uline{un schéma quasi}-\uline{compact}, \uline{ayant un faisceau ample}, \uline{et} @@M12@@ \uline{un morphisme propre et parfait}. \uline{On suppose de plus vérifiée l'une des deux conditions suivantes}:
@@M13@@
\uline{de façon que le} «\uline{théorème de finitude}» \uline{soit vrai} (III 2.2). \uline{On se donne un anneau} @@M14@@, \uline{et} \uline{on considère l'homomorphisme} «\uline{image directe}» (IV 2.11):
@@M15@@
\uline{On suppose enfin qu'il existe un élément} @@M16@@ \uline{tel que} @@M17@@ \uline{soit inversible}. \uline{Alors l'homomorphisme} :
@@M18@@
\uline{est injectif}.

On écrit la formule de projection:
@@M19@@
Or @@M20@@ est d'augmentation inversible, donc est inversible (6.3). Par suite, l'homomorphisme composé :
